% ===========================================================================
%  sviluppatore/04-motore/06-confronto.tex
% ===========================================================================
\section{Confrontare un live con il suo backtest}\label{sez:live-confronto}

La prova che chiude il cerchio: una strategia intraday gira per qualche
giorno su un conto demo, poi si rifà il backtest sui dati del live e si
confrontano le due equity. Perché le curve coincidano servono le cose già
viste --- il motore che esegue e poi sorveglia (\vediSez{barra-trigger}), lo
stream che riceve ciò che il backtest ha (\vediSez{live-stream}) --- e due
cose che riguardano il backtest: che arrivi alla fine dei dati, e che giri
sulle stesse barre del live.

\subsection{Il walk-forward arriva alla fine dei dati}

\figuraTikz{fold-di-coda}
  {Le finestre del walk-forward con le frazioni di default. L'ultimo fold
   regolare si chiude al \perc{95} del periodo; il fold di coda ha la stessa
   ampiezza e finisce alla fine dei dati, ma nell'equity OOS entra solo il
   tratto che nessun altro copriva.}
  {fold-di-coda}

\begin{description}
  \item[La fine vuota] \cod{end\_date} vuota vuol dire «fino agli ultimi
    dati» e si risolve a \emph{domani}
    (\simbolo{script\_config\_helper.\_resolve\_end\_date}): i caricatori
    tagliano a quell'istante, e con «oggi» le barre intraday di oggi
    restavano fuori. La fine effettiva la decide l'ultima barra comune ai
    simboli (\simbolo{run\_walkforward.\_compute\_effective\_range}). Prima
    era oggi meno sette giorni, «per il ritardo di yfinance»: la settimana
    tolta escludeva sempre i dati più recenti, cioè i giorni del live.
  \item[L'ora] Per un timeframe intraday una fine senza ora diventa
    \cod{23:59:59} di quel giorno, e la fine di ogni fold conserva l'ora
    (\simbolo{\_istante\_str}: la data sola se è mezzanotte, come sempre per
    il giornaliero). Prima un fold intraday si fermava alla mezzanotte del
    suo ultimo giorno.
  \item[Il benchmark] accorcia la fine solo se finisce più di tre giorni
    prima dei simboli operativi (\cod{BENCHMARK\_RITARDO\_MAX}): è un
    download vecchio. SPY intraday chiude alle 21 UTC e i cambi alle 22, e
    ogni sera il backtest di una strategia sui cambi perdeva l'ultimo tratto.
  \item[Il fold di coda] Con le frazioni di default l'ultimo \perc{5} del
    periodo non entrava in nessun test: su due mesi di barre M15 sono gli
    ultimi tre giorni, proprio quelli di un live appena concluso.
    \simbolo{\_fold\_di\_coda} aggiunge un fold la cui finestra di test
    finisce alla fine dei dati; lo segna \cod{coda = True} sull'istanza
    (\simbolo{e\_coda}), così i checkpoint vecchi restano leggibili. Non
    conta fra i fold delle statistiche: numero, dispersione degli Sharpe,
    idoneità (\file{validation/robustness.py}) restano quelli dei fold
    regolari.
  \item[Le operazioni OOS] seguono l'equity: ogni fold porta solo quelle fra
    la fine del precedente (esclusa) e la propria. Le finestre di test si
    sovrappongono, e quelle del tratto in comune comparivano due volte.
    L'ordinamento dei rendimenti concatenati è stabile
    (\cod{kind="mergesort"}), così a parità di istante «il primo fold» è
    davvero il più vecchio.
  \item[I checkpoint] passano alla versione \cod{"2"} e ricordano anche le
    frazioni: un checkpoint fatto con altre finestre non si riusa.
\end{description}

\subsection{Il backtest sulle barre del live}\label{sez:barre-live}

Dal vivo una strategia decide sulle barre del broker (\vediSez{live-contesto}),
che con il salvataggio dei dati del live acceso restano in
\file{data/live/\_contesto/<fonte>/<TF>/<SIMBOLO>.parquet} (per MT5 la fonte
è \cod{mt5\_<server>}). Il backtest legge il deposito, scaricato da un
provider, e anche dopo l'unione il periodo del live vi arriva dal provider
del file: un file non mescola provenienze (\vediSez{live-dati-vivo}). Sono
serie diverse. Sulle M15 dei cambi i rendimenti di Yahoo e di Dukascopy
hanno correlazione fra \num{0.95} e \num{0.97}; su un live simulato di due giorni
il backtest sulle barre Yahoo faceva \num{13} operazioni contro le \num{21}
del live, \num{9} in comune.

\begin{center}
\begin{tabular}{@{}L{46mm}L{92mm}@{}}
\toprule
\textbf{Pezzo} & \textbf{Che cosa fa} \\
\midrule
\rigaTenue
\file{data/barre\_live.py}
  & \simbolo{fonti} elenca le fonti con barre salvate e i loro timeframe;
    \simbolo{per\_strategia} dice quali hanno, al timeframe primario della
    strategia, tutti i suoi simboli operativi (con periodo e numero di
    barre); \simbolo{per\_simboli} dà le barre per simbolo e timeframe, e un
    timeframe aggiuntivo che la fonte non ha lo ricava dalle sue barre del
    primario. \\
\file{dashboard/ui\_ext\_barre\_live.py}
  & \cod{GET /api/data/live\_bars}: le fonti e, per strategia, le fonti
    utilizzabili. \\
\rigaTenue
\file{runner.js}, \file{v2-run.js}
  & la voce \voce{Run on the live bars} nel menu della scheda, con
    timeframe, fonte e periodo; una conferma; poi \cod{POST
    /api/run/strategy} con \cod{live\_bars}. \\
\simbolo{api\_run\_strategy}
  & controlla che la fonte esista e lancia il job con
    \cod{QS\_BARRE\_LIVE=<fonte>} nell'ambiente. \\
\rigaTenue
\file{execution/run\_walkforward.py}
  & con \cod{QS\_BARRE\_LIVE}, i simboli che la fonte ha si leggono dalle
    sue barre (righe \cod{[BARS]} nel log) e gli altri dal deposito ---
    benchmark compreso, se il broker non lo quota; nessuna barra per
    nessun simbolo è un errore. Il deposito non cambia; il risultato lo
    annota in \cod{config.bars\_source}. \\
\bottomrule
\end{tabular}
\end{center}

\noindent
Per guardare solo la finestra del live, le metriche hanno l'equity barra per
barra (\cod{equity\_bars}, \cod{returns\_bars}, \cod{bars\_timeframe}:
\vediSez{nuova-metrica}): l'esempio 6 di \file{metrics/METRIC\_TEMPLATE.md}
ribasa l'equity del backtest dall'ora d'avvio del live al capitale del conto.

\subsection{La prova completa}

\file{tests/prova\_live\_intraday.py} prova i pezzi uno per uno. La prova
intera usa quattro strategie di prova (\cod{live\_test\_rules}, regole fisse
su M15 e H1; \cod{live\_test\_model}, una regressione logistica
riaddestrata ogni quattro ore dentro la strategia; \cod{live\_test\_dollar},
DXY e VIX orari da Yahoo con la sua porta dei dati; \cod{live\_test\_all},
le tre insieme), un broker MT5 finto che dà le barre Dukascopy fino a un
orologio simulato e le quotazioni minuto per minuto lungo la barra, un
provider Yahoo finto, la porta dei dati con Yahoo che pubblica dieci minuti
dopo l'ora, tutto nell'ambito del live (\file{data.vivo}) e un conto
\simbolo{PaperBroker}: \num{200} cicli M15 (due giorni) più la sentinella
ogni minuto, poi l'unione dei dati del live e il backtest sulle barre del
live.

\begin{center}
\begin{tabular}{@{}lrrrr@{}}
\toprule
\textbf{Strategia} & \textbf{Operazioni} & \textbf{Resa backtest}
  & \textbf{Scarto stream} & \textbf{Scarto conto} \\
\midrule
\rigaTenue
\cod{live\_test\_rules}  & \num{21}/\num{21} & \perc{\num{+0.044}} & \perc{0.0003} & \perc{0.013} \\
\cod{live\_test\_model}  & \num{90}/\num{90} & \perc{\num{+0.030}} & \perc{0.0013} & \perc{0.042} \\
\rigaTenue
\cod{live\_test\_dollar} & \num{44}/\num{44} & \perc{\num{-0.211}} & \perc{0.0001} & \perc{0.033} \\
\cod{live\_test\_all}    & \num{64}/\num{64} & \perc{\num{-0.098}} & \perc{0.0001} & \perc{0.043} \\
\bottomrule
\end{tabular}
\end{center}

\noindent
Operazioni: quelle del live che il backtest ritrova, sul totale. Scarto: il
massimo, sul periodo del live, fra l'equity del backtest e quella dello
stream o del conto, ribasate all'inizio. Il conto esegue ai prezzi del minuto
in cui il ciclo gira, non all'apertura della barra: è la differenza che resta
anche su un conto vero, insieme ai costi del broker.
